\chapter{Kết luận}

\section{Tổng kết}

Khóa luận này xuất phát từ một quan sát đơn giản về cách các hệ thống điều phối xe đi chung
đang vận hành: chúng báo cho hành khách một giờ hẹn, rồi tiếp tục tối ưu lại và dời giờ ấy nhiều
lần, trong khi không một thước đo chuẩn nào của ngành ghi nhận điều đó. Câu hỏi đặt ra là liệu
một hệ thống có thể đưa ra lời hứa mà nó thực sự giữ được hay không, và cái đó đắt bao nhiêu.

Để trả lời, chúng tôi xây dựng cơ chế lời hứa có biên được chứng nhận, với ba thành phần. Thành
phần thứ nhất là phép phân rã ba mốc, dựng một mốc phản thực bằng cách chiếu lại chính kế hoạch
cũ từ vị trí thật hiện tại của xe, nhờ đó tách được phần dịch chuyển do nhà điều hành gây ra
khỏi phần thay đổi khi giữ nguyên kế hoạch. Thành phần thứ hai là sổ cái cam kết riêng cho từng hành
khách, cộng dồn theo thời gian và đối chiếu với một trần, được kiểm tra ngay tại thời điểm ra
quyết định. Thành phần thứ ba là hệ thống chứng cứ niêm phong, bảo đảm rằng mọi so sánh giữa
các cấu hình đều kiểm chứng được tới từng lần chạy.

Về mặt lý thuyết, đóng góp trung tâm là khung so sánh bốn họ ràng buộc lời hứa và Định lý
\ref{dl:hocong}: xét riêng phần cam kết, chỉ ràng buộc theo hành vi không bao giờ là nguyên nhân
làm rỗng tập phương án khi thế giới trôi dạt, còn họ nào trong các họ khác --- hạn chót một chiều
kiểu Schulz và Pfeiffer, hạn chót hai chiều, ràng buộc theo kết quả khách nhìn thấy --- hỏng thì
phụ thuộc chiều trôi. Nền của khung này là ba định lý về hai cơ sở tính sổ cái. Tính theo hành vi
của nhà điều hành thì không bao giờ bị ràng buộc cam kết làm rỗng, vì phương án không làm gì theo
định nghĩa không gây ra gì. Tính theo kết quả mà hành khách nhìn thấy thì có thể khiến hệ thống
không còn bất kỳ quyết định hợp lệ nào, vì ngay cả việc không làm gì cũng bị tính phí trôi dạt.
Và khi một điều kiện (H) đúng ở mọi lần công bố, chính sách tối ưu theo cơ sở kết quả không phục
vụ nhiều hơn chính sách tối ưu theo cơ sở hành vi; trên panel, (H) hỏng ở 228 lần công bố đã chọn
của B1, nên đây là một mệnh đề có điều kiện chứ không phải kết luận về panel. Các chứng minh đều
sơ cấp.
Phân biệt hai cơ sở trùng với phân biệt cổ điển giữa nghĩa vụ phương tiện và nghĩa vụ kết quả trong
luật hợp đồng, và điểm mới nằm ở chỗ mốc tham chiếu ngoại sinh \emph{tính được} ranh giới ``ngoài
tầm kiểm soát'' mà luật xử lý bằng văn xuôi.

Về mặt thực nghiệm, 558 lần chạy trên một panel 20 ô đã niêm phong cho bốn nhóm kết quả. Nhóm
thứ nhất là quy trách: trong chế độ không cam kết, phần do bộ điều phối tự gây ra bằng
$89{,}06\%$ lượng dịch chuyển mà hành khách nhìn thấy, và thiệt hại dồn vào một cái đuôi mỏng --- $23{,}2\%$ số
hành khách hứng toàn bộ, người chịu thiệt nặng nhất bị dời gần mười bốn phút mà hệ thống vẫn
tính là phục vụ thành công. Nhóm thứ hai là đường cong giá của cam kết trên chín cấu hình, với
phát hiện rằng phần đắt nhất là quyết định cam kết chút nào đó chứ không phải độ chặt của lời
hứa về thời gian. Nhóm thứ ba nhất quán với các định lý, kể cả dưới dữ liệu tắc đường đo thật của
Manhattan (thăm dò), nơi cơ sở kết quả vô nghiệm do cam kết ở 19 trên 20 ô còn cơ sở hành vi không
vô nghiệm do cam kết ở ô nào --- dù 11 ô vẫn dừng vì bất biến cửa sổ đón của tầng miền. Một thí
nghiệm thăm dò có tiền đăng ký sau đó đặt hai họ hạn chót vào cùng phép thử: dự đoán theo từng ô,
niêm phong trước khi chạy các nhánh được dự đoán, đúng ở $98/102$ cặp trên panel chính và ở $31/32$, $40/40$ cặp trên khung
sáng và khung tối của một tập giữ lại. Buổi sáng, hạn chót 30 giây tính từ lời hứa đầu vô nghiệm do
cam kết ở 19/20 ô; buổi tối, khi tắc đường giảm, hạn chót một chiều vô nghiệm ở 0/8 ô trong khi hạn
chót hai chiều và cổng kết quả vô nghiệm ở 7/8 và 8/8 ô; cổng hành vi không vô nghiệm do cam kết ở
lần chạy nào trong 56 lần. Mức 30 giây chặt hơn nhiều so với 15 phút của Schulz và Pfeiffer. Khi không
có trôi dạt, hai luật cùng 30 giây nhận đúng cùng một tập khách ở 20/20 ô. Vì vậy đây là kết quả về
vùng lời hứa chặt dưới tắc đường. Một phép quét thăm dò từ 30 giây tới 15 phút (mục \ref{muc:quet})
cho thấy hai luật phục vụ gần như cùng số khách ở mọi mức. Khác biệt là lời hứa có giữ được hay
không. Dưới tắc đường tự tạo có cú sốc, ở mức chặt, lời hứa của hạn chót bị phá với $10$--$19\%$ khách,
còn trần hành vi giữ mọi lời hứa. Trên ngày tắc đường đo thật, hạn chót chỉ phá lời hứa ở mức 30 giây
buổi sáng, với $2{,}2\%$ khách. Nhóm thứ tư kiểm tra độ vững: chín thước đo cho cùng một thứ hạng, không
một ngoại lệ.

Ngoài ra, khóa luận ghi nhận bốn kết quả phụ mà chúng tôi cho là đáng chú ý. Thứ nhất, sáu cơ
chế bảo vệ giờ đón khác nhau --- kể cả cơ chế bảo đảm một chiều mà văn liệu sử dụng --- rơi vào
đúng hai hành vi, vì hình thức của một lời bảo đảm chỉ có ý nghĩa trong chừng mực bộ sinh
phương án có thể thực thi nó. Thứ hai, giá của một lời hứa không phải hằng số mà tăng theo mức
tự do của nhà điều hành, và mặt cong hai chiều ở mục \ref{muc:matcong} cho thấy hai yếu tố này
\emph{tương tác}: trần càng chặt thì việc cho thêm tự do càng đắt, từ $0{,}42$ lên $1{,}44$
điểm phần trăm khi trần siết từ 300 xuống 30 giây. Thứ ba, hệ quả của định lý về hai cơ sở
biểu hiện theo kiểu phân loại chứ không theo kiểu biên --- cơ sở kết quả mất nguyên cả ô chứ
không phục vụ ít hơn vài hành khách --- và lý do là số học: khả thi cấp ô sụp theo $(1-p)^{n}$
với $n$ là số hành khách mỗi ô.

Thứ tư, ba mức cải tiến bộ sinh phương án mà chúng tôi dự định thực hiện, cùng bộ sinh biết sắp
xếp lại, đều được đánh giá bằng chính dữ liệu của hệ thống trước khi viết mã. Cắt tỉa sớm là phép
toán rỗng vì chưa từng có ứng viên khả thi nào bị bỏ hay nhánh nào chưa khai triển. Xếp hạng vị trí
chèn cũng thực chất rỗng: trần số ứng viên chưa bao giờ chạm nên xếp hạng không đổi được đầu vào
của bộ giải, và trên các kỷ nguyên có nhận hoặc từ chối khách chỉ $2{,}48\%$ tập của một xe có từ
hai vị trí chèn trở lên. Đặt trước ngân sách chạy được nhưng trùng khớp tới từng quyết định với đối
chứng; tầng ngay trên nó là một hằng số nên không che nó, còn tầng số khách nhận ở trên cùng thì có
thể đã phân định mọi so sánh --- nguyên nhân chưa xác định được từ bằng chứng hiện có. Và bộ sinh
biết sắp xếp lại, vốn đã có sẵn, va vào một bất biến thứ hai của mô hình về thời điểm một chặng hết
sửa được. Điểm chung của cả bốn là chúng được đánh giá
bằng phép đo rẻ tiền trên bằng chứng đã có, chứ không bằng nhiều ngày chạy mô phỏng --- và bằng
chứng ấy do chính cơ chế lời hứa sinh ra.

\section{Hạn chế}

Chúng tôi muốn nêu các hạn chế một cách rõ ràng, vì một số trong đó ảnh hưởng trực tiếp tới
cách đọc kết quả.

Hạn chế lớn nhất là \textbf{quy mô}. Panel 108 yêu cầu và 8 xe mỗi ô nhỏ hơn nhiều so với các
thực nghiệm quy mô Manhattan trong văn liệu. Lựa chọn này có chủ ý --- mục tiêu là so sánh có
kiểm soát giữa các cấu hình trên cùng đầu vào --- nhưng hệ quả là các con số tuyệt đối không so
trực tiếp được với văn liệu.

Hạn chế thứ hai là \textbf{phần trùng lặp với công trình gần nhất}. Như mục \ref{muc:doichieu} đã trình bày chi
tiết, sáu trong số các hạng mục của khóa luận đã có trong Fielbaum và Alonso-Mora (2020): khái
niệm, đo độ lớn, đếm số lần, phân loại nguyên nhân định tính, đường cong đánh đổi, và cơ chế
ratchet. Chúng tôi coi phần đó là nền chứ không phải đóng góp, và đã cố gắng ghi rõ ranh giới ở
bảng \ref{bang:trung}. Hai công trình mới hơn còn làm hẹp thêm phạm vi: Elmi và cộng sự đã có bảo
đảm lời hứa theo từng khách, và Schulz và Pfeiffer đã có một hạn chót cưỡng chế được tính từ lời
hứa đầu (mục \ref{muc:chodung}).

Hạn chế thứ ba là \textbf{trôi dạt ở panel chính chỉ là cận dưới}. Bộ dữ liệu thử gốc có đồ thị
thời gian di chuyển tĩnh, nên trôi dạt đo được đến từ độ trễ thực thi chứ không từ tắc đường.
Thí nghiệm bơm trôi dạt đã bù một phần nhưng chưa trọn vẹn: nó cho kết quả về tính khả thi nhưng
không cho kết quả về tỉ lệ phục vụ, vì một bất biến miền của hệ thống chặn lại. Hệ số trôi dạt
được áp vô hướng cho toàn mạng, và các lần chạy có trôi dạt chưa qua bộ kiểm chứng độc lập của
dự án, nên kết quả trôi dạt chỉ là thăm dò. ``Mạng'' của bộ dữ liệu thử cũng không phải mạng đường
phố mà là bao đóng thời gian đi giữa 96 điểm đã chọn.

Hạn chế thứ tư là \textbf{mức bằng chứng cho khung bốn họ cổng}. Định lý \ref{dl:hocong} được
chứng minh, và phần đối chiếu với hệ thật là một thí nghiệm có tiền đăng ký, nhưng vẫn là
\emph{thăm dò}: các lần chạy có trôi dạt không qua bộ kiểm chứng của dự án, trôi dạt là một hệ số vô
hướng, mỗi khung của tập giữ lại chỉ có 8 ô của hai ngày và một hạt giống, và quy tắc dự đoán là
heuristic --- năm ca sai đều là ca nhánh được thử rời xa quỹ đạo tham chiếu. Chiều trôi thuận lợi
mới được thấy trên một loại khung (buổi tối của hai ngày). Ngoài ra định lý chỉ nói về phần cam kết: dưới trôi dạt thật,
hệ thống dùng trần hành vi vẫn dừng ở 11 trên 20 ô vì bất biến cửa sổ đón, và định lý không che
được giới hạn ấy.

Hạn chế thứ năm là \textbf{lập luận phúc lợi kém chắc chắn hơn phần kỹ thuật}. Nó dựa trên hệ số
định giá từ một bối cảnh địa lý khác và trên giả định rằng mọi chuyến bị mất có cùng thặng dư.
Chúng tôi trình bày nó như một ngưỡng chứ không phải một kết luận.

Cuối cùng, khóa luận \textbf{không tuyên bố vượt trội về tỉ lệ phục vụ}. Mọi cấu hình có cam
kết đều phục vụ ít hành khách hơn cấu hình không cam kết. Đóng góp nằm ở việc biến một đại
lượng trước nay chỉ đo được sau sự việc thành một ràng buộc cưỡng chế được trước khi quyết
định, và ở việc chỉ ra loại ràng buộc lời hứa nào có thể tự nó làm hệ thống kẹt.

\section{Hướng phát triển}

Từ những hạn chế trên, chúng tôi đề xuất bảy hướng tiếp theo, xếp theo thứ tự ưu tiên mà chúng
tôi cho là hợp lý.

\textbf{Đưa thí nghiệm bốn họ cổng lên mức chính thức và mở rộng thế giới trôi dạt.} Thí nghiệm ở
mục \ref{muc:tang2} là thăm dò. Bước tiếp theo là cho bộ kiểm chứng của dự án chấp nhận các lần chạy
có trôi dạt, để kết quả thành chính thức; rồi thay hệ số vô hướng bằng trôi dạt theo từng cạnh, gồm
cả trôi hai chiều và hồi phục sau cú sốc, trên nhiều ngày và nhiều hạt giống hơn. Khi đó mới đo được
dự đoán phụ chưa quan sát được --- hạn chót một chiều phục vụ nhiều khách hơn trần hành vi khi tắc
đường giảm; trong phép quét ở mục \ref{muc:quet}, hai luật phục vụ bằng nhau ở thế giới tắc đường
giảm, nhưng ở đó gần như mọi khách đều được phục vụ nên phép thử chưa đủ nhạy --- và mới biết quy tắc dự đoán theo ô còn đúng không khi nhánh được thử rời xa quỹ đạo
tham chiếu.

\textbf{Sổ cái chỉ tính phần gây hại.} Cơ sở hành vi phạt cả một sửa chữa có lợi (phản ví dụ CE-3
ở bảng \ref{bang:phanvidu}). Một biến thể tự nhiên chỉ tính phần quyết định làm lời hứa \emph{xa
lời hứa đầu hơn} so với giữ nguyên tuyến,
\[
h_t = \max\big(0,\; |p_t - p_0| - |e_t - p_0|\big).
\]
Theo bất đẳng thức tam giác $0 \le h_t \le A_t$, và phương án không làm gì có $h_t = 0$, nên phần
(a) của Định lý \ref{dl:hocong} vẫn giữ nguyên; CE-3, nếu đó là lần công bố đầu sau lời hứa đầu
($p_0 = v = 600$), cho $h = \max(0, 0 - 120) = 0$ và được miễn. Nhưng trên hệ hiện tại khác biệt
gần như bằng không: trên B1, tổng $h$ là $64.821{,}2$ giây so với tổng $A$ là $64.863{,}1$ giây ---
chênh $41{,}9$ giây ($0{,}06\%$), cùng 403 khách bị tính --- dù $h \ne A$ ở 158 trên 479 lần công
bố. Vì bộ sinh chèn thuần luôn cho $p \ge e$, $h \ne A$ đúng khi tuyến giữ nguyên đã sớm hơn lời
hứa đầu ($e < p_0$): khi ấy phần quyết định dời lời hứa về phía $p_0$ không bị tính. Không
có sửa chữa kiểu CE-3, tức kéo lời hứa sớm hơn $e$, nào để miễn; biến thể này chỉ có giá trị rõ khi
có bộ sinh biết sắp xếp lại. Chúng tôi chưa rà văn liệu cho riêng biến thể này, nên không tuyên bố
gì về tính mới của nó.

\textbf{So tỉ lệ phục vụ giữa hai cơ sở trên một panel có ít hành khách hơn mỗi ô.} Định lý
\ref{dl:tapcon} là một mệnh đề về chính sách tối ưu, không đo trực tiếp được; điều đo được là tỉ lệ
phục vụ của bộ điều phối đang cài. Mục \ref{muc:troidat} cho thấy phép so
sánh tỉ lệ phục vụ giữa hai cơ sở không lấy được trên panel hiện tại, vì khả thi cấp ô sụp theo
$(1-p)^{n}$ với $n \approx 80$. Lối thoát tự nhiên là giảm $n$: một panel nhiều ô nhỏ hơn, mỗi
ô vài chục hành khách, sẽ giữ được nhiều ô khả thi ở cùng mức trôi dạt, nhờ đó có tập ô chung
để so. Đây là thay đổi về thiết kế panel chứ không phải về hệ thống, nên chi phí thấp.

\textbf{Mở rộng mặt cong sang chiều tự do thứ ba.} Mục \ref{muc:matcong} đã dựng mặt cong trên hai mức tự
do và cho thấy tự do với độ chặt \emph{tương tác} chứ không cộng độc lập. Chiều tự do ở đây là
cửa sổ đón; hai chiều tự nhiên tiếp theo là cỡ đội xe và mật độ nhu cầu. Nếu quy luật giữ nguyên
trên cả ba chiều thì có thể rút ra một quy tắc chọn trần theo mức tự do, thay vì phải đo lại cho
từng cấu hình vận hành.

\textbf{Cửa sổ đón thương lượng có chọn lọc.} Như mục \ref{muc:dathu} đã nêu, hướng này chưa bị bác bỏ. Bài
toán đã được đặt đúng với cả trần lẫn kiểu hỏng đều đã đo: lấy được bao nhiêu trong 172 hành
khách mà vẫn giữ chờ trung vị quanh 407 giây. Tuy nhiên nó đòi hỏi sửa mã nguồn và mở rộng phát
biểu vấn đề sang kiểm soát tiếp nhận, nên cần cân nhắc kỹ về phạm vi.

\textbf{Nới lỏng bất biến miền về cửa sổ đón.} Phát hiện ở mục \ref{muc:batbien} --- rằng hệ thống chưa biểu
diễn được trạng thái mà riêng giao thông phá vỡ một cửa sổ đã cam kết --- là một giới hạn thật
của mô hình. Cơ chế nới lỏng theo mốc ngoại sinh đã tồn tại trong hệ thống cho bộ sinh phương án;
mở rộng nó sang tầng kiểm tra trạng thái sẽ cho phép nghiên cứu chế độ cao điểm một cách đầy đủ.

\textbf{Đóng băng chặng đang chạy, rồi mới mở khóa bộ sinh biết sắp xếp lại.} Mục
\ref{muc:bosinh} cho thấy hệ thống đã có sẵn cơ chế sắp xếp lại nhưng nó hỏng trên 18 trên 20
ô, vì mô hình miền chỉ đóng băng một điểm dừng khi xe \emph{tới nơi}, còn bộ mô phỏng khóa chặng
ngay khi xe \emph{lăn bánh}. Việc cần làm có phạm vi rõ ràng và phải theo đúng thứ tự: thêm
chuyển trạng thái đóng băng chặng đang bay ở tầng miền --- kiểu dữ liệu lộ trình đã cho phép
trạng thái ``đã đóng băng nhưng chưa thực hiện'', chỉ thiếu điểm gọi; cho chính sách vector cứng
nhận hạn mức sửa chữa; và nới điều kiện kiểm tra cấu hình để hai thứ đi được cùng nhau. Chỉ sau
ba bước đó mới đặt được câu hỏi mà hướng này thật sự muốn hỏi: liệu quyền sắp xếp lại có thu hồi
được những hành khách mà cổng cam kết đang từ chối hay không. Chúng tôi lưu ý rằng ô duy nhất
chạy trọn vẹn cho kết quả bằng đúng nhánh không cam kết, nên kỳ vọng nên đặt thấp.

\section{Kết luận cuối}

Đóng góp mà chúng tôi tin là bền nhất của khóa luận không phải một con số cụ thể, mà là một cách
đặt vấn đề. Thước đo kinh điển của văn liệu tham chiếu tới giá trị thực tế, chỉ biết được sau khi
hành khách đã được đón, nên về mặt toán học không thể trở thành ràng buộc; các công trình gần đây
đã kiểm soát hiện tượng bằng những ràng buộc trên chính lời hứa --- cấm tăng cho một nhóm khách và
giới hạn trong độ chùng cho những khách còn lại, hoặc một hạn chót tính từ lời hứa đầu. Khóa luận
này chặn một đối tượng khác: một trần có mức chỉ trên phần do nhà điều hành gây ra, tách khỏi phần
thay đổi khi giữ nguyên kế hoạch, định nghĩa trên các đại lượng đều biết được tại thời điểm quyết
định.

Và khi biến một thước đo thành một ràng buộc, một câu hỏi mới xuất hiện: \emph{tính sổ cái theo
cái gì}. Câu trả lời quyết định cổng cam kết có thể tự nó làm hệ thống kẹt hay không. Định lý
\ref{dl:hocong} phát biểu điều đó cho bốn họ ràng buộc, và một thí nghiệm có tiền đăng ký trên dữ
liệu tắc đường thật xác nhận nó ở mức thăm dò, theo từng ô và ở cả hai chiều trôi: ràng buộc theo kết
quả khách nhìn thấy --- cách tự nhiên nhất để biến thước đo của văn liệu thành ràng buộc --- kẹt ở
cả hai chiều; một hạn chót một chiều tính từ lời hứa đầu kẹt khi tắc đường tăng và an toàn khi tắc
đường giảm; hạn chót hai chiều kẹt ở cả hai, ở cận trên buổi sáng và ở cận dưới buổi tối; chỉ trần
hành vi thì không bao giờ là nguyên nhân.
